\begin{figure}[!hbtp]
    \centering
    \resizebox{0.9\columnwidth}{!}{%
    \begin{tikzpicture}[
        % Estilos herdados do seu padrão
        recntagle/.style={
            rectangle,
            rounded corners=3pt,
            inner sep=6pt,
            text centered,
            font=\small\bfseries,
            draw=black!70,
            fill=blue!10,
            line width=0.8pt,
            text width=4.6cm,
            minimum height=1.2cm
        },
        arrow/.style = {
            thick, 
            ->, 
            >={Stealth[length=2.5mm]}
        },
        sub_block/.style = {
            rectangle,
            rounded corners=3pt,
            inner sep=5pt,
            text centered,
            draw=black!70,
            fill=blue!5,
            line width=0.7pt,
            font=\small\bfseries,
            text width=5.8cm,
            minimum height=1.6cm
        }
    ]

        % ==========================================
        % Definição das Questões de Pesquisa
        % ==========================================
        \node (definicao) [recntagle] {Definição das Questões de Pesquisa};

        % Subblocos: questões de pesquisa
        \node (qp1) [sub_block] at ([xshift=7.2cm, yshift=1.2cm]definicao.center) {QP1: Quais são os fundamentos teóricos e arquiteturais do \textit{Blockchain}?};
        \node (qp2) [sub_block] at ([xshift=7.2cm, yshift=-1.2cm]definicao.center) {QP2: Como os requisitos dos setores financeiro e de saúde impulsionaram a evolução do \textit{Blockchain}?};

        % Setas do bloco de definição para as questões de pesquisa
        \draw [arrow] (definicao.east) -- ++(0.9cm,0) |- (qp1.west);
        \draw [arrow] (definicao.east) -- ++(0.9cm,0) |- (qp2.west);

        % ==========================================
        % Fluxo de pesquisa
        % ==========================================
        \node (busca) [recntagle, below=1.4cm of definicao] {Processo de Busca};
        \node (criterios) [recntagle, below=1.4cm of busca] {Critérios de Inclusão e Exclusão de Estudos};

        % Setas do fluxo principal
        \draw [arrow] (definicao.south) -- (busca.north);
        \draw [arrow] (busca.south) -- (criterios.north);

    \end{tikzpicture}%
    }
    \caption{Fluxo metodológico da pesquisa: a definição das questões de pesquisa (QP1 e QP2) orienta o processo de busca, seguido pela aplicação dos critérios de inclusão e exclusão de estudos.}
    \label{fig:fluxo_metodologico}
\end{figure}
